% ======================================================================
% Section 6 -- Results II: Explaining performance (~1.5 pp).
% Direct register; no em dashes; no "ladder/rung/stem" vocabulary.
% BINDING: multi-seed error bars; parameter/budget matching; linear
% axes; uniform y-span; one star per fixed k; menu at true cost.
% ======================================================================
\section{Results II: Explaining Performance}
\label{sec:performance}

The selection procedure of \cref{sec:framework-game2} runs identically
on all \unfrozen{\ResLeadPanels} datasets and in two model families,
GCCN message-passing routes and HOPSE encoding blocks. We compare
against each family's own architecture sweep, with every method priced
by its total training compute in FLOPs under one measurement
convention (\cref{app:flops-conventions}). All markers carry
multi-seed error bars over identical seeds. \Cref{tab:performance-summary}
(\cref{app:selection-table}) tabulates every marker with exact values,
the $B{=}1$ arm, and per-row cost multiples. Comparisons against
published fixed architectures and fixed-$k$ variants are in
\cref{app:selection}.

\begin{figure}[!t]
  \centering
  \IfFileExists{../figures/out/lead_xai_anchored.pdf}{%
    \includegraphics[width=\linewidth]{../figures/out/lead_xai_anchored.pdf}%
  }{%
    \figureplaceholder{figures/out/lead_xai_anchored.pdf}%
  }
  \caption{\textbf{Architecture selection on the GraphXAI datasets.}
  Test accuracy vs.\ total training compute (FLOPs, linear axes).
  Dark markers are the architecture chosen by our selection procedure
  in each family; pale markers are each family's sweep at its full
  grid cost. Bars are sd over seeds.}
  \label{fig:lead}
\end{figure}

\begin{figure}[!t]
  \centering
  \IfFileExists{../figures/out/lead_bench_anchored.pdf}{%
    \includegraphics[width=\linewidth]{../figures/out/lead_bench_anchored.pdf}%
  }{%
    \figureplaceholder{figures/out/lead_bench_anchored.pdf}%
  }
  \caption{\textbf{Architecture selection on the TopoBench
  benchmarks.} Same conventions as \cref{fig:lead}. Top row:
  complex-level tasks; bottom row: node-level citation complexes. The
  HOPSE family is evaluated on complex-level tasks, so the node-level
  panels show the GCCN family. On MUTAG and NCI109 the GCCN sweep
  marker is the best of the 64-configuration retraining enumeration
  at estimated cost. MANTRA is scored in balanced accuracy (majority
  class \unfrozen{\ResMantraMajorityRate}).}
  \label{fig:lead-bench}
\end{figure}

\expclaim{Selection concedes at most \unfrozen{\ResSelMaxSweepGap}
accuracy for a
\unfrozen{\ResSweepCostMultMin}--\unfrozen{\ResSweepCostMultMax}$\times$
compute reduction.} In \cref{tab:performance-summary}, the best
selection arm matches or exceeds the best sweep within one sweep
seed-sd on \unfrozen{\ResSelWithinRows} of
\unfrozen{\ResSelTotalRows} datasets. On the remaining datasets the
sweep ends higher, never by more than \unfrozen{\ResSelMaxSweepGap}.
The largest concessions sit where the masked game carries the least
signal: the node-level citation complexes, whose small validation
splits raise the noise floor the stop rule reads (next claim), the
class-imbalanced MANTRA, and Alkane-Carbonyl, where the comparator is
the exhaustive enumeration rather than a sweep. The
$B{=}2$ budget costs
$\sim$\unfrozen{\ResLadderBTwoCostMult}$\times$ more, keeps a
\unfrozen{\ResSweepVsBTwoMin}--\unfrozen{\ResSweepVsBTwoMax}$\times$
compute margin over the sweep, and stabilizes the one dataset where
$B{=}1$ is noisy (Fluoride-Carbonyl,
\cref{tab:performance-summary}). The same procedure, unchanged, runs
in the HOPSE family (\cref{fig:lead,fig:lead-bench}). These
costs use the masked value; evaluating the retraining value instead
costs \unfrozen{\ResSweepEqMin}--\unfrozen{\ResSweepEqMax}
sweep-equivalents, so the masked value carries every cost claim
(regret $\le$ \unfrozen{\ResRegretMax} against retraining-value
selections, \cref{app:accounting}).

\expclaim{The number of neighborhoods kept is not a hyperparameter.}
The stop derives from the validation noise floor $\sqrt{p(1-p)/m}$,
which matches measured retraining noise in order of magnitude on all
\unfrozen{\ResBinomOrderDatasets} datasets. The rule was frozen before
its final scoring. Out of sample, the sequence of candidates contains
the exact optimum on
\unfrozen{\ResLadderOOSContained}/\unfrozen{\ResLadderOOSTotal}
never-touched architecture landscapes, with worst regret
\unfrozen{\ResLadderOOSWorstRegret}. Fixing $k$ in advance is the
alternative, and it fails where the right size differs: on MUTAG the
fixed $k{=}3$ model collapses while the selection holds near the
exact optimum (\cref{tab:performance-summary}).

\expclaim{Cross-family gaps are model-class facts that no
within-family selection can close.} On Benzene, HOPSE's spectral
encodings clear the ceiling of the entire GCCN neighborhood space by
$\sim$\unfrozen{\ResBenzeneSpectralMargin}. On MANTRA orientability,
the GCCN family sits near the count-only score of
\unfrozen{\ResCountingCeiling} in balanced accuracy while the same
procedure in the HOPSE family clears it at
$\sim$\unfrozen{\ResMantraSweepVsLadderMult}$\times$ less compute
than the HOPSE sweep (\cref{fig:lead-bench}). The GCCN sweep makes
the same point at full grid cost: its validation-best configuration
is a majority collapse (\cref{tab:performance-summary}). Plain-accuracy MANTRA
numbers collapse to the majority class, so all MANTRA results are
balanced accuracy (\cref{app:experimental}).

